\section{Thảo luận và hạn chế}\label{sec:discussion}

\paragraph{Vì sao các ca còn lọt qua.} Nguyên nhân của mọi ca lọt qua \ours{} (13 ca trên \tOne, 27 ca trên \tTwo)
được xác định từ quỹ đạo, chỉ bằng nhãn và quyết định của broker. Mỗi ca rơi vào đúng một chỗ hở của
Mệnh đề~\ref{prop:scope}. Mọi ca \mSL{} và \mAL{} (9 và 16 ca) vi phạm giả định (iii). Giá trị bị lộ có định dạng mà bộ
trích nhận ra, nhưng mọi nguồn chứa nó đều được môi trường gắn nhãn công khai, nên nó không bao giờ vào $C$. Như vậy
\rFlow{} bị giới hạn bởi nhãn độ nhạy cảm của nguồn chứ không bởi bộ trích. Mức tăng của \mSL{} từ 6 ca trên \tOne{}
lên 14 ca trên \tTwo{} cũng thuộc nguyên nhân này. Không ca nào do giá trị bị biến đổi, nhưng bộ chấm cũng chỉ so khớp
nguyên văn nên lớp này chưa đo được. Ba ca \mUS, một ca \mPT{} của \tOne{} và ca \mCR{} duy nhất vi phạm giả định (ii).
Bốn ca đầu gửi tới một địa chỉ có trong $u$ nhưng không được cấp quyền. Ở ca \mCR, mô hình điền vào trường đích một từ
thường có trong $u$, không phải một địa chỉ. Mười ca \mPT{} còn lại (3 và 7) thuộc $H_\mathit{tamper}$, nằm ngoài mệnh
đề. Chỉ \rSig{} chặn được lớp này, và chỉ khi bộ phát hiện gắn cờ đoạn chèn. Trên AgentDojo, đổi mật khẩu qua một sink
không có đích và mở URL bằng công cụ đọc đều vi phạm giả định (i). Tập thích nghi ở Mục~\ref{sec:adaptive} không nhắm
vào các chỗ hở này. \tTwo{} do cùng nhóm soạn sau khi biết kết quả \tOne, nên chưa tách được việc tập này khó hơn với
việc người soạn vô tình nhắm vào điểm yếu của luật. Một bộ tấn công hộp trắng sinh từ chính các giả định trên là bước
đánh giá cần làm tiếp.

\paragraph{Ngôn ngữ.} Ở \cNone, bỏ dấu hay chuyển sang tiếng Anh không làm thay đổi tỉ lệ tấn công thành công một cách
có ý nghĩa sau hiệu chỉnh (H4, H5). Khoảng tin cậy vẫn chứa chênh lệch tới khoảng 12 điểm, nên bài báo không khẳng định
tiếng Việt dễ hay khó bị tấn công hơn. Yếu tố chịu ảnh hưởng rõ của ngôn ngữ là phòng thủ dựa trên cụm từ. Trên Dev, bộ
phát hiện cụm từ bật ở 14 đến 19 trên 48 tấn công tiếng Việt chuẩn nhưng chỉ ở 3 trên 48 tấn công tiếng Anh. Trên
AgentDojo, \rSig{} không bật lần nào ở tiếng Anh nhưng bật 14 lần ở Slack tiếng Việt. \rDest{} và \rFlow{} không phụ
thuộc ngôn ngữ, còn \rSig{} thì có.

\paragraph{Hạn chế.} Các tập niêm phong do cùng một nhóm soạn, dù người soạn không đọc mã phòng thủ. Dữ liệu là tổng
hợp, công cụ là bản giả và nhiệm vụ chỉ có một lượt. Các mô hình mở chỉ có 8 đến 12 tỉ tham số. \ours{} làm giảm tiện
ích lành tính khoảng 5 điểm. Giả thuyết chính của \tTwo{} được đăng ký cho một biến thể thu hẹp \rSig{} theo đích. Biến
thể này tăng tiện ích 1,0 điểm với cùng tỉ lệ tấn công thành công, nhưng khoảng tin cậy theo họ chứa 0, nên nó không
được đưa vào bài. Gemini 3 Flash chỉ chạy năm cấu hình trên Dev và cấu hình không phòng thủ trên AgentDojo. Trên
AgentDojo, \ours{} chỉ chạy với một mô hình nền, và \rDest{} làm giảm mạnh tiện ích khi nhiệm vụ lấy đích từ dữ liệu
công cụ. Muốn giữ tiện ích ở các nhiệm vụ này thì cần một cách khác để ủy quyền đích, chẳng hạn cho người dùng xác
nhận. Tham chiếu chặn mọi sink là ước lượng từ quỹ đạo. Phản hồi từ chối chưa được tách riêng. Độ phủ của
bộ trích trên toàn bộ quỹ đạo và độ trễ chưa được đo.
